\section{Smart Contract Implementation}

The smart contract layer forms the core of the SkillChain platform, handling business logic, fund management, and access control. Three interconnected contracts work together to provide complete functionality.

\subsection{UserProfile Contract}

\subsubsection{Purpose and Functionality}

The UserProfile contract manages user registration, profiles, and role assignment. It ensures each user has a unique identity on the platform and maintains the mapping between Ethereum addresses and user profiles.

\subsubsection{Data Structures}

\begin{lstlisting}[language=Solidity, caption={UserProfile Data Structures}]
enum Role { 
    Client,      // 0 - Can post jobs
    Freelancer   // 1 - Can submit proposals
}

struct Profile {
    uint256 userID;        // Auto-incremented unique ID
    string username;       // Unique identifier
    string cid;           // IPFS Content Identifier
    Role role;            // Client or Freelancer
    address userAddress;  // Ethereum wallet address
}
\end{lstlisting}

\subsubsection{Key Functions}

\begin{table}[H]
\centering
\caption{UserProfile Functions}
\begin{tabular}{|p{4cm}|p{2.5cm}|p{4cm}|p{2cm}|}
\hline
\textbf{Function} & \textbf{Visibility} & \textbf{Description} & \textbf{Gas Cost} \\
\hline
setProfile & external & Creates/updates profile with validation & $\sim$65,000 \\
\hline
getProfileByAddress & external view & Retrieves profile by wallet & $\sim$2,500 \\
\hline
getUserByUsername & external view & Reverse lookup by username & $\sim$2,000 \\
\hline
\end{tabular}
\end{table}

\subsubsection{Profile Creation Algorithm}

The profile creation process follows these steps:

\begin{enumerate}
    \item Validate username and CID are non-empty.
    \item Increment user ID counter.
    \item Check username uniqueness (either new profile or changed username).
    \item Store profile data in mapping.
    \item Update username ownership mapping.
    \item Emit ProfileSet event for indexing.
\end{enumerate}

\subsection{JobsContract}

\subsubsection{Purpose and Functionality}

The JobsContract handles job creation, escrow management, and payment release. It ensures funds are securely held and only released upon mutual agreement.

\subsubsection{Data Structures}

\begin{lstlisting}[language=Solidity, caption={JobsContract Data Structures}]
struct Job {
    uint256 jobID;              // Unique job identifier
    string jobCID;              // IPFS reference to job details
    string deliverableCID;      // IPFS reference to completed work
    bool completed;             // Job completion status
    address client;             // Job poster
    address freelancer;         // Hired freelancer
    uint256 amount;             // Escrowed payment (wei)
    bool clientApproved;        // Client approval flag
    bool freelancerApproved;    // Freelancer approval flag
}
\end{lstlisting}

\subsubsection{Key Functions}

\begin{table}[H]
\centering
\caption{JobsContract Functions}
\begin{tabular}{|p{4cm}|p{2.5cm}|p{6cm}|}
\hline
\textbf{Function} & \textbf{Modifier} & \textbf{Description} \\
\hline
CreateJob & payable & Creates job with ETH escrow deposit \\
\hline
HireFreelancer & external & Assigns freelancer to job \\
\hline
MarkFLJobCompleted & external & Freelancer marks job complete \\
\hline
MarkClientJobCompleted & external & Client marks job complete \\
\hline
\_releaseFunds & private & Releases payment to freelancer \\
\hline
\end{tabular}
\end{table}

\subsubsection{Job Creation Algorithm}

\begin{enumerate}
    \item Validate msg.value is greater than 0 (funds required).
    \item Increment job ID counter.
    \item Create Job struct with provided CID and deposited funds.
    \item Store job in \_jobs mapping.
    \item Append job ID to client's job list.
    \item Emit JobCreated event.
\end{enumerate}

\subsubsection{Payment Release Algorithm}

The payment release function implements critical security patterns:

\begin{lstlisting}[language=Solidity, caption={Payment Release Function}]
function _releaseFunds(uint256 jobID) private {
    Job storage job = _jobs[jobID];
    if (job.freelancerApproved && job.clientApproved && !job.completed) {
        // Effects: Update state first (prevent re-entrancy)
        job.completed = true;
        uint256 payment = job.amount;
        job.amount = 0;  // Critical: Set to zero before external call
        
        // Interactions: External call last
        (bool sent, ) = job.freelancer.call{value: payment}("");
        require(sent, "Transfer failed");
        
        emit JobCompleted(msg.sender, jobID);
    }
}
\end{lstlisting}

\textbf{Security Pattern:} The contract uses Checks-Effects-Interactions pattern with re-entrancy protection by setting job.amount = 0 before external calls.

\subsection{ProposalsContract}

\subsubsection{Purpose and Functionality}

The ProposalsContract manages proposal submissions and approval workflow. It integrates with JobsContract to automatically hire freelancers upon proposal approval.

\subsubsection{Interface Definition}

\begin{lstlisting}[language=Solidity, caption={Jobs Contract Interface}]
interface IJobsContract {
    function getJobOwner(uint256 jobID) external view returns (address);
    function HireFreelancer(address freelancerAddress, uint256 jobID) external;
}
\end{lstlisting}

\subsubsection{Data Structures}

\begin{lstlisting}[language=Solidity, caption={ProposalsContract Data Structures}]
struct Proposal {
    uint256 proposalID;     // Unique proposal ID
    string proposalCID;     // IPFS reference to proposal details
    uint256 jobID;          // Reference to target job
    bool approved;          // Approval status
    address freelancer;     // Proposal submitter
}
\end{lstlisting}

\subsubsection{Proposal Approval Algorithm}

\begin{enumerate}
    \item Retrieve proposal from storage.
    \item Verify proposal is not already approved.
    \item Query JobsContract to verify job ownership.
    \item Ensure caller is the job owner.
    \item Mark proposal as approved.
    \item Call JobsContract.HireFreelancer to assign freelancer.
    \item Emit ProposalApproved event.
\end{enumerate}

\section{Frontend Implementation}

\subsection{Project Structure}

The frontend follows a well-organized directory structure:

\begin{lstlisting}[language=bash, caption={Frontend Directory Structure}]
client/
|-- src/
|   |-- app/                    # Next.js App Router pages
|   |   |-- layout.tsx         # Root layout with providers
|   |   |-- page.tsx           # Landing page
|   |   |-- dashboard/         # User dashboard
|   |   |-- explore/jobs/      # Job listings
|   |   |-- job/[id]/          # Job details
|   |   |-- track/job/[id]/    # Client tracking
|   |   |-- deliver/job/[id]/  # Freelancer deliverables
|   |   |-- sign-up/           # Profile creation
|   |-- components/
|   |   |-- ui/               # UI primitives (Button, Input, etc.)
|   |   |-- navbar.tsx        # Navigation component
|   |   |-- job-card.tsx      # Job listing cards
|   |   |-- popups/           # Dialog and modal components
|   |-- providers/
|   |   |-- provider.tsx      # Wagmi + RainbowKit configuration
|   |-- store/
|   |   |-- user.store.ts     # Zustand user state
|   |   |-- transition.store.ts # UI transition state
|   |-- lib/
|   |   |-- pinata.ts         # IPFS operations
|   |   |-- utils.ts          # Utility functions
|   |-- types/
|   |   |-- user.types.ts     # TypeScript interfaces
|   |   |-- job.types.ts
|   |-- abi/                  # Contract ABIs
\end{lstlisting}

\subsection{State Management}

\subsubsection{Zustand User Store}

\begin{lstlisting}[language=TypeScript, caption={User State Interface}]
interface UserState {
  user: IUser | null;
  isConnected: boolean;
  isRegistered: boolean;
}

interface UserActions {
  setUser: (user: IUser) => void;
  clearUser: () => void;
  setConnected: (status: boolean) => void;
}
\end{lstlisting}

State persists in localStorage for seamless reconnection across sessions.

\subsection{IPFS Integration}

\subsubsection{Upload Function}

\begin{lstlisting}[language=TypeScript, caption={IPFS Upload Function}]
async function uploadJSONToPinata(obj: any): Promise<string> {
  const blob = new Blob([JSON.stringify(obj)], { 
    type: "application/json" 
  });
  
  const file = new File([blob], `${obj.username}.json`, { 
    type: "application/json" 
  });
  
  const formData = new FormData();
  formData.append("file", file);
  
  const res = await fetch("https://api.pinata.cloud/pinning/pinFileToIPFS", {
    method: "POST",
    headers: {
      Authorization: `Bearer ${PINATA_JWT}`,
    },
    body: formData,
  });
  
  const data = await res.json();
  return data.IpfsHash;
}
\end{lstlisting}

\subsubsection{Fetch Function}

\begin{lstlisting}[language=TypeScript, caption={IPFS Fetch Function}]
async function fetchFromPinata(cid: string): Promise<any> {
  const res = await fetch(
    `https://${PINATA_GATEWAY}/ipfs/${cid}`
  );
  return await res.json();
}
\end{lstlisting}

\subsection{Blockchain Configuration}

\begin{lstlisting}[language=TypeScript, caption={Wagmi Configuration}]
const config = getDefaultConfig({
  appName: "SkillChain",
  projectId: "YOUR_PROJECT_ID",
  chains: [foundry],
  transports: {
    [foundry.id]: http("http://127.0.0.1:8545"),
  },
});
\end{lstlisting}

\section{Algorithms and Data Structures}

\subsection{Data Structure Analysis}

\begin{table}[H]
\centering
\caption{Data Structure Complexity Analysis}
\begin{tabular}{|p{3.5cm}|p{2.5cm}|p{3cm}|p{3cm}|p{3cm}|}
\hline
\textbf{Structure} & \textbf{Type} & \textbf{Time Complexity} & \textbf{Space Complexity} & \textbf{Purpose} \\
\hline
\_profiles & Mapping & O(1) access & O(n) & User lookup by address \\
\hline
\_usernameOwner & Mapping & O(1) access & O(n) & Reverse username lookup \\
\hline
\_jobs & Mapping & O(1) access & O(n) & Job lookup by ID \\
\hline
\_clientJobIDs & Mapping + Array & O(1) + O(n) & O(n) & Client job history \\
\hline
\_jobProposalIDs & Mapping + Array & O(1) + O(n) & O(n) & Job proposals list \\
\hline
\_proposals & Mapping & O(1) access & O(n) & Proposal lookup \\
\hline
\_freelancerProposalIDs & Mapping + Array & O(1) + O(n) & O(n) & Freelancer proposals \\
\hline
\end{tabular}
\end{table}

\subsection{Complete Job Lifecycle Algorithm}

\begin{enumerate}
    \item \textbf{Phase 1: Setup}
    \begin{itemize}
        \item User A registers as Client via UserProfile.setProfile()
        \item User B registers as Freelancer via UserProfile.setProfile()
    \end{itemize}
    
    \item \textbf{Phase 2: Job Creation}
    \begin{itemize}
        \item Client calls JobsContract.CreateJob\{value: X\}(jobCID)
        \item Contract stores job with escrow funds
        \item Emit JobCreated event
    \end{itemize}
    
    \item \textbf{Phase 3: Proposal Submission}
    \begin{itemize}
        \item Freelancer calls ProposalsContract.CreateProposal(propCID, jobID)
        \item Contract stores proposal reference
        \item Emit ProposalCreated event
    \end{itemize}
    
    \item \textbf{Phase 4: Proposal Approval}
    \begin{itemize}
        \item Client calls ProposalsContract.ApproveProposal(proposalID)
        \item Contract verifies job ownership
        \item Contract calls JobsContract.HireFreelancer(freelancer, jobID)
        \item Emit ProposalApproved event
    \end{itemize}
    
    \item \textbf{Phase 5: Work Execution}
    \begin{itemize}
        \item Freelancer completes work off-chain
    \end{itemize}
    
    \item \textbf{Phase 6: Completion Marking}
    \begin{itemize}
        \item Freelancer calls JobsContract.MarkFLJobCompleted(jobID, deliverableCID)
        \item Contract sets freelancerApproved = true
        \item If clientApproved already true, call \_releaseFunds(jobID)
        \item Client calls JobsContract.MarkClientJobCompleted(jobID)
        \item Contract sets clientApproved = true
        \item If freelancerApproved already true, call \_releaseFunds(jobID)
    \end{itemize}
    
    \item \textbf{Phase 7: Payment Release}
    \begin{itemize}
        \item \_releaseFunds() verifies both approvals
        \item Transfers escrow to freelancer
        \item Marks job completed
        \item Emit JobCompleted event
    \end{itemize}
\end{enumerate}

\subsection{Gas Optimization Techniques}

The smart contracts implement several gas optimization strategies:

\begin{enumerate}
    \item \textbf{Calldata vs Memory:} Using calldata for function parameters reduces gas costs by avoiding memory allocation.
    \item \textbf{Storage Packing:} Struct fields are ordered to minimize storage slots (each slot is 32 bytes).
    \item \textbf{View Functions:} Read-only operations marked as view cost no gas when called externally.
    \item \textbf{Event Logging:} Expensive storage operations replaced with events where appropriate.
    \item \textbf{Early Returns:} Validation checks at function start prevent wasted computation.
\end{enumerate}

\subsection{Access Control Matrix}

\begin{table}[H]
\centering
\caption{Function Access Control}
\begin{tabular}{|p{5cm}|p{4cm}|p{5cm}|}
\hline
\textbf{Function} & \textbf{Access Control} & \textbf{Validation} \\
\hline
setProfile & Any address & Username uniqueness \\
\hline
CreateJob & Any address & msg.value $>$ 0 \\
\hline
HireFreelancer & Client or ProposalsContract & Job exists, not completed \\
\hline
CreateProposal & Any address & Job exists \\
\hline
ApproveProposal & Job owner only & Proposal exists, not approved \\
\hline
MarkFLJobCompleted & Assigned freelancer only & Job not completed \\
\hline
MarkClientJobCompleted & Job client only & Job not completed \\
\hline
\end{tabular}
\end{table}